\section{Introduction}\label{sec:introduction}


Evaluating the literary quality of narrative fiction represents a fundamental challenge at the intersection of computational analysis and humanistic interpretation. As modern literary production increasingly intersects with developments in painting, music, and other art forms, a broader convergence with generative artificial intelligence is becoming inevitable. This convergence manifests at varying degrees: while recent years have seen the emergence of LLM-generated novels, images, and short films, such works have also attracted substantial criticism. In many cases, a clear gap remains between these generated artifacts and human-authored works of recognized literary quality, particularly with respect to depth, coherence, and interpretive richness. Consequently, the central problem is not merely how to generate creative content, but how to evaluate it in a principled and meaningful way.

At present, both human and computational approaches to literary evaluation face limitations. Readers, critics, and scholars routinely engage in nuanced assessments of narrative works, yet such judgments are difficult to formalize, scale, or reproduce. Existing computational approaches to text quality assessment primarily focus on surface-level properties such as grammatical correctness, lexical diversity, and local coherence, often operationalized through n-gram overlap or embedding-based similarity metrics such as BLEU, ROUGE, and BERTScore \cite{papineni2002bleu,lin2004rouge,zhang2019bertscore}. While useful for specific applications, these metrics have been shown to correlate weakly with human judgments and to inadequately capture higher-level narrative and interpretive qualities \cite{novikova2017we,reiter2018structured}. As a result, they fail to address the interpretive, experiential, and normative dimensions that play a central role in literary criticism.

This mismatch reflects a deeper conceptual issue rather than a purely technical limitation. Literary evaluation does not correspond to a single, theory-neutral notion of quality that can be reduced to a scalar score. Instead, it encompasses multiple semi-independent dimensions operating at different epistemic levels. Some dimensions concern observable textual and structural properties that can be analyzed through reproducible linguistic and narrative methods. Other dimensions involve interpretive constructions (cultural representation, emotional–psychological experience, and philosophical meaning) that emerge from interactions between text, reader, and critical framework. This layered view of narrative assessment is consistent with prior work in computational narratology and story evaluation, which emphasizes the multidimensional nature of coherence, engagement, and meaning in narrative texts \cite{graesser2004coh,callan2023interesting,yang2024makes}. Recognizing this structure suggests that effective evaluation frameworks must be ontologically grounded, explicitly distinguishing descriptive analysis from interpretive modeling and normative valuation. Without such structure, evaluation risks conflating fundamentally different types of judgments: reproducible textual measurements may be confused with theory-dependent interpretations, and normative valuations may be treated as objective properties of texts.

Building on this observation, we propose SAGE, a framework that adopts ontology as a methodological construct for structuring literary evaluation rather than as a claim about objective aesthetic value. The framework organizes literary assessment into six analytically distinct layers (see Table~\ref{tab:framework_layers}), each corresponding to a different class of evaluative entities and epistemic commitments. Lower layers (L1–L3) assess observable textual properties through rule-based computational metrics, while higher layers (L4–L6) evaluate interpretive constructs, namely cultural representation, emotional–psychological experience, and existential–philosophical engagement, through theory-grounded LLM-based analysis. These layers are non-substitutable: higher-level judgments presuppose but do not collapse into lower-level features, and evaluation at each layer is constrained to its own ontological domain. By making explicit the evaluative commitments underlying each dimension, the framework enables controlled study of evaluation behavior without conflating reproducible textual measurements with theory-dependent interpretations.

Our empirical investigation examines evaluation patterns across narratives that differ in provenance and institutional reception: canonical literary works, commercial genre fiction, and LLM-generated stories. These categories are not treated as absolute indicators of literary quality, but as contrasting contexts in which different evaluative frameworks have historically been applied. We introduce two evaluation conditions. In content-limit evaluation, assessment is based directly on the textual content of a literary work. In title-limit evaluation, assessment relies on external critical knowledge (the work's title or reputation), approximating expert judgment without direct textual access. Comparing these conditions allows us to analyze how different sources of information shape literary evaluation and to probe the grounding of interpretive judgments, consistent with prior benchmarks for human and automatic story evaluation \cite{chhun2022human,hamilton2025narrabench}.

Our study is guided by four research questions:

\textbf{RQ1 (Reliability):} Can LLM-based literary evaluation produce stable and reproducible assessments across multiple instances and independent evaluators within explicitly defined evaluative layers?

\textbf{RQ2 (Discriminative Capacity):} Do LLM evaluations successfully distinguish between texts of different literary provenance (canonical, popular, LLM-generated), demonstrating meaningful quality differentiation rather than arbitrary scoring?

\textbf{RQ3 (Evaluation Robustness):} How do different evaluation conditions (direct textual analysis versus reliance on critical context) affect LLM assessments, and when do evaluations remain stable and interpretable?

\textbf{RQ4 (Dimensional Independence):} Do different evaluative layers exhibit empirical independence, supporting the theoretical claim that literary evaluation consists of distinct but complementary dimensions?

Through systematic experiments addressing these questions, we demonstrate that ontology-aware LLM-based evaluation can yield stable and interpretable assessments within specified interpretive layers, while also revealing limitations in grounding evaluative judgments without adequate contextual or theoretical support. By making explicit the assumptions and commitments underlying different dimensions of literary evaluation, this work contributes a principled foundation for computational engagement with literary texts.